\section{A constructive boundary: positive higher-order dispersion}
\label{sec:superorder}

A positive addition of strictly higher degree has the opposite
dilation sign and admits negative-energy stationary waves arbitrarily
close to the lower-order critical mass. This comparison lies outside
both the upper dilation condition and the original order-$m$ growth
class. The proof uses standard constrained minimization. Its
fourth/second-order and fractional-order special cases have substantial
precedent; it is included to explain the variational mechanism in
arbitrary orders, not as a counterexample to the unmodified homogeneous
equation or as an isolated sharpness test of the upper dilation condition.

Fix $N\ge1$, $k>m>1$, and $\varepsilon>0$. Set
\[
 r=\frac mN,\qquad p=2+2r,\qquad
 A_j(f)=\int |\xi|^j|\widehat f(\xi)|^2\,d\xi,
\]
and consider
\begin{equation}\label{higher:equation}
 iu_t=(|D|^m+\varepsilon|D|^k)u-|u|^{2m/N}u,
 \qquad
 E_\varepsilon(f)=\frac12A_m(f)+\frac\varepsilon2A_k(f)
                    -\frac1p\|f\|_p^p.
\end{equation}
Let $M_{c,m}$ be the sharp critical mass for the \emph{lower-order}
symbol $|\xi|^m$:
\begin{equation}\label{higher:gn}
 \|f\|_p^p\le\frac p2
       \left(\frac{M(f)}{M_{c,m}}\right)^r A_m(f).
\end{equation}
Define
\[
 I_\varepsilon(M)=\inf\{E_\varepsilon(f):
      f\in H^{k/2}(\mathbb R^N;\mathbb C),\ M(f)=M\}.
\]

\begin{theorem}[No gap above the lower-order threshold]
\label{higher:existence}
For $0<M\le M_{c,m}$, $I_\varepsilon(M)=0$ and the infimum is not
attained. For every $M>M_{c,m}$,
\[
 -\infty<I_\varepsilon(M)<0,
\]
and every minimizing sequence is precompact in $H^{k/2}$ after
translations and passage to a subsequence. There is a real minimizer
$Q_M\in L^\infty\cap H^{k+1}$ satisfying
\begin{equation}\label{higher:profile}
 (|D|^m+\varepsilon|D|^k)Q_M+\omega_MQ_M
       -|Q_M|^{2m/N}Q_M=0,
 \qquad \omega_M>0.
\end{equation}
Its energy obeys
\begin{equation}\label{higher:energy}
 E_\varepsilon(Q_M)
       =-\frac{k-m}{2m}\varepsilon A_k(Q_M)<0.
\end{equation}
Thus $e^{i\omega_Mt}Q_M$ is a global stationary solution of
\eqref{higher:equation}, with zero momentum and strictly negative
energy, for every $M>M_{c,m}$.
\end{theorem}

\begin{proof}
\emph{Coercivity and sign.}
Highest-order Gagliardo--Nirenberg interpolation gives
\begin{equation}\label{higher:coercive-gn}
 \|f\|_p^p\le C_{N,m,k}
 M(f)^{1+m/N-m/k}A_k(f)^{m/k}.
\end{equation}
Since $m/k<1$, Young's inequality implies, at fixed mass,
\begin{equation}\label{higher:coercivity}
 E_\varepsilon(f)\ge\frac\varepsilon4A_k(f)
                         -C_{N,m,k,\varepsilon,M}.
\end{equation}
The infimum is finite and minimizing sequences are bounded in
$H^{k/2}$. If $M\le M_{c,m}$, \eqref{higher:gn} gives
\[
 E_\varepsilon(f)\ge
 \frac12\left[1-(M/M_{c,m})^r\right]A_m(f)
                       +\frac\varepsilon2A_k(f)>0
\]
for every nonzero admissible $f$. On the other hand, for any smooth
datum of mass $M$, the mass-preserving dilation
$f_\lambda(x)=\lambda^{N/2}f(\lambda x)$ satisfies
\begin{equation}\label{higher:dilation}
 E_\varepsilon(f_\lambda)
       =\lambda^mE_0(f)+\frac\varepsilon2\lambda^kA_k(f)
       \longrightarrow0\quad(\lambda\downarrow0).
\end{equation}
This proves the first assertion. If $M>M_{c,m}$, multiply a
lower-order sharp optimizer by $\sqrt{M/M_{c,m}}$ to obtain negative
$E_0$. Approximation in $H^{m/2}$ by real $C_c^\infty$ functions,
followed by a mass adjustment, gives such an $f$ of exact mass $M$.
Equation~\eqref{higher:dilation} is negative for sufficiently small
$\lambda$, because $k>m$. This does not require the lower-order
optimizer itself to belong to $H^{k/2}$.

\emph{Strict binding.}
For $\theta>1$, amplitude scaling gives
\[
 E_\varepsilon(\sqrt\theta f)
 =\theta E_\varepsilon(f)
       -\frac{\theta^{p/2}-\theta}{p}\|f\|_p^p
 \le\theta E_\varepsilon(f),
\]
and also
\[
 E_\varepsilon(\sqrt\theta f)
 =\theta^{p/2}E_\varepsilon(f)
  +\frac{\theta-\theta^{p/2}}2(A_m(f)+\varepsilon A_k(f))
 \le\theta^{p/2}E_\varepsilon(f).
\]
Consequently $I_\varepsilon(M)/M$ is nonincreasing, with
\begin{equation}\label{higher:strict-amplitude}
 I_\varepsilon(\theta M)\le\theta^{p/2}I_\varepsilon(M)
               <\theta I_\varepsilon(M)
 \quad\text{if }I_\varepsilon(M)<0.
\end{equation}
For $M>M_{c,m}$ and $M=M_1+M_2$ with $M_1,M_2>0$, it follows that
\begin{equation}\label{higher:binding}
 I_\varepsilon(M)<I_\varepsilon(M_1)+I_\varepsilon(M_2).
\end{equation}
Indeed, if both pieces are subthreshold the right side is zero.
Otherwise the monotonicity of $I_\varepsilon(M)/M$, strict for a
piece with negative infimum by~\eqref{higher:strict-amplitude}, proves
\eqref{higher:binding} upon multiplying by the two masses and adding.

\emph{Compactness and attainment.}
Let $f_n$ be a minimizing sequence at $M>M_{c,m}$ and put $s=k/2$.
Local mass vanishing is impossible. To see this, a smooth unit-cube
partition and local GN interpolation give
\[
 \|\chi_j f_n\|_{2+2k/N}^{2+2k/N}
    \le C\|\chi_j f_n\|_2^{2k/N}\|\chi_j f_n\|_{H^s}^2,
 \qquad
 \sum_j\|\chi_j f_n\|_{H^s}^2\le C\|f_n\|_{H^s}^2.
\]
Thus $\sup_y\int_{B(y,1)}|f_n|^2\to0$ would imply
$\|f_n\|_{2+2k/N}\to0$, and interpolation would imply
$\|f_n\|_p\to0$ since $2<p<2+2k/N$. This contradicts the negative
limiting energy.

After translations and a subsequence, there is therefore a weak
$H^s$ limit $Q\ne0$, with local strong $L^2$ convergence and
almost-everywhere convergence. Write $w_n=f_n-Q$,
$M_1=M(Q)>0$, and $M_2=M-M_1\ge0$. Hilbert splitting of both
kinetic forms and the Brezis--Lieb lemma yield
\begin{align}
 M(f_n)&=M(Q)+M(w_n)+o(1),\notag\\
 E_\varepsilon(f_n)&=E_\varepsilon(Q)
                         +E_\varepsilon(w_n)+o(1).
 \label{higher:splitting}
\end{align}
If $M_2>0$, normalize $w_n$ to exactly that mass. Its normalization
factor tends to one and all energy terms remain bounded, so the
energy changes by $o(1)$. Equation~\eqref{higher:splitting} then
contradicts~\eqref{higher:binding}. Hence $M_1=M$ and $M(w_n)\to0$.
Equation~\eqref{higher:coercive-gn} now gives $\|w_n\|_p\to0$,
so~\eqref{higher:splitting} forces
\[
 E_\varepsilon(Q)=I_\varepsilon(M),\qquad
 A_m(w_n)+\varepsilon A_k(w_n)\longrightarrow0.
\]
Thus convergence is strong in $H^{k/2}$.

A real minimizer can be obtained without a diamagnetic inequality.
The real and complex infima coincide: write $f=A+iB$ and, if both
component masses are positive, put
\[
 \vartheta=M(A)/M(f),\qquad
 g=A/\sqrt\vartheta,\qquad h=B/\sqrt{1-\vartheta}.
\]
Both $g,h$ are real and have mass $M(f)$. The two real-even kinetic
forms split, while convexity gives
\[
 |f|^p=\big(\vartheta g^2+(1-\vartheta)h^2\big)^{p/2}
             \le\vartheta|g|^p+(1-\vartheta)|h|^p.
\]
Therefore $E_\varepsilon(f)\ge
\vartheta E_\varepsilon(g)+(1-\vartheta)E_\varepsilon(h)$,
which is at least the real infimum. The cases of a zero component
are immediate. A real minimizing sequence and the preceding
compactness proof give a real minimizer $Q_M$.

\emph{Stationary identities.}
The energy is $C^1$ on $H^{k/2}$. The Lagrange multiplier theorem
gives~\eqref{higher:profile} weakly for a real $\omega_M$.
Minimizing along the mass-preserving dilation gives directly
\begin{equation}\label{higher:pohozaev}
 \frac m2A_m(Q_M)+\frac{k\varepsilon}2A_k(Q_M)
                         -\frac mp\|Q_M\|_p^p=0.
\end{equation}
Only the explicit scalar formula~\eqref{higher:dilation} is
differentiated; no spatial moment or energy-space dilation generator
is needed. Equation~\eqref{higher:energy} follows immediately.
Testing the weak equation against $Q_M$ additionally gives
\begin{equation}\label{higher:frequency}
 \omega_M M=\frac mN A_m(Q_M)
   +\left(\frac{k}{m}+\frac{k}{N}-1\right)
           \varepsilon A_k(Q_M)>0.
\end{equation}

\emph{Regularity.}
For $\omega_M>0$, the resolvent
\[
 R=(|D|^m+\varepsilon|D|^k+\omega_M)^{-1}
\]
maps $L^q$ to $H^{k,q}$ for every $1<q<\infty$. Its multiplier
after multiplication by $(1+|\xi|^2)^{k/2}$ satisfies the Mihlin
bounds, including at low frequency after each derivative is weighted
by the corresponding power of $|\xi|$. The equation reads
$Q_M=R(|Q_M|^{2m/N}Q_M)$.
If $N<k$, the starting energy embedding gives $L^\infty$; if $N=k$,
arbitrarily large finite input exponents give the same conclusion.
If $N>k$, start at
\[
 \beta_0=\frac12-\frac{k}{2N},\qquad
 \beta_{j+1}=(1+2m/N)\beta_j-k/N,
\]
where $\beta_j$ is the reciprocal integrability exponent. The fixed
point is $k/(2m)>1/2>\beta_0$, so finitely many iterations give an
input exponent above $N/k$ and hence $L^\infty$. The first nonlinear
input exponent exceeds one because $2N/(N-k)>1+2m/N$; all later
ones improve. An endpoint can be bypassed by slightly weakening an
intermediate exponent. Boundedness gives the nonlinearity in $L^2$,
so $Q_M\in H^k\subset H^1$. The $C^1$ chain rule then puts the
nonlinearity in $H^1$, and one more resolvent step gives
$Q_M\in H^{k+1}$. This proves the theorem.
\end{proof}

\subsection{Quantitative degeneration at the lower-order threshold}

\begin{proposition}[Low-frequency bounds]\label{higher:low-frequency}
For the minimizers in Theorem~\ref{higher:existence}, set
\[d=(M/M_{c,m})^{m/N}-1>0.\]
Then
\begin{align}
 \varepsilon A_k(Q_M)&\le\frac mk d A_m(Q_M),
       \label{higher:ratio}\\
 A_k(Q_M)&\le M\left(\frac{m d}{k\varepsilon}\right)^{k/(k-m)},
 & A_m(Q_M)&\le M\left(\frac{m d}{k\varepsilon}\right)^{m/(k-m)},
       \label{higher:kinetic-bounds}\\
 -\frac{k-m}{2k}M d
  \left(\frac{m d}{k\varepsilon}\right)^{m/(k-m)}
     &\le I_\varepsilon(M)<0.
       \label{higher:energy-bound}
\end{align}
Thus $A_m$, $A_k$, $|E_\varepsilon|$, and $\omega_M$ tend to zero
as $M\downarrow M_{c,m}$ with $\varepsilon$ fixed. After normalizing
$A_m=1$, these minimizers converge subsequentially in $H^{m/2}$,
modulo translations and constant phases, to a lower-order sharp
optimizer.
\end{proposition}
\begin{proof}
Combining~\eqref{higher:gn} and~\eqref{higher:pohozaev} gives
\eqref{higher:ratio}. Fourier interpolation gives
$A_m\le M^{1-m/k}A_k^{m/k}$, which proves
\eqref{higher:kinetic-bounds}; the energy identity gives
\eqref{higher:energy-bound}. Equation~\eqref{higher:frequency}
then gives $\omega_M\to0$.

For $M_n\downarrow M_{c,m}$, put
\[
 \lambda_n=A_m(Q_{M_n})^{-1/m},\qquad
 \widetilde Q_n(x)=\lambda_n^{N/2}Q_{M_n}(\lambda_nx).
\]
These profiles satisfy $A_m(\widetilde Q_n)=1$, $M(\widetilde Q_n)=M_n$,
and
\[
 E_0(\widetilde Q_n)
 =-\frac{k\varepsilon A_k(Q_{M_n})}{2mA_m(Q_{M_n})}<0,
 \qquad |E_0(\widetilde Q_n)|\le d_n/2\longrightarrow0.
\]
The normalized sharp-GN compactness lemma therefore gives the
asserted strong convergence. The higher-order coefficient in the
normalized equation is
\[
 \varepsilon\lambda_n^{m-k}
       =\varepsilon A_m(Q_{M_n})^{(k-m)/m}\longrightarrow0.
\]
Thus this term becomes small but still determines the negative
energy in~\eqref{higher:energy}.
\end{proof}

If a lower-order optimizer $Q$ additionally belongs to $H^{k/2}$,
one obtains matching energy powers by using the trial family
$\sqrt{M/M_{c,m}}\,Q_\lambda$. Put $a=A_m(Q)$, $b=A_k(Q)$,
and $\vartheta=M/M_{c,m}$, with $M(Q)=M_{c,m}$. Its optimal scale
is $\lambda^{k-m}=m a d/(k\varepsilon b)$, yielding
\begin{equation}\label{higher:energy-upper}
 I_\varepsilon(M)\le
 -\frac{\vartheta(k-m)}{2k}a d
       \left(\frac{m a d}{k\varepsilon b}\right)^{m/(k-m)}.
\end{equation}
Together with~\eqref{higher:energy-bound}, this gives the same power
$\varepsilon^{-m/(k-m)}d^{k/(k-m)}$ on both sides, not an exact
asymptotic coefficient. The additional optimizer regularity is
retained here for arbitrarily large $k$; the existence theorem does
not require it.

\begin{remark}[Meaning of the boundary]
For $b(\xi)=\varepsilon|\xi|^k$ one has
$\xi\cdot\nabla b=kb>mb$ away from zero. This is precisely the
opposite dilation sign and accounts for~\eqref{higher:energy}.
Every positive coefficient admits negative-energy stationary waves
at masses arbitrarily close to $M_{c,m}$. This modification also violates
the order-$m$ growth bound, so by itself it does not isolate the role of
the upper dilation condition within that growth class. A small
coefficient here is not a small bounded quadratic-form perturbation on
$H^{m/2}$, and the modified symbol is not homogeneous of degree $m$;
the fixed-order homogeneous robustness theorem is unaffected.

For a strong $H^{k/2}$ flow with conservation and the usual norm
continuation criterion,~\eqref{higher:coercivity} additionally rules
out finite-time divergence of the energy norm. This conditional
flow consequence is separate from the explicitly global stationary
solutions proved above.
\end{remark}
